\subsection{Definition} \label{IV.1}
The chapter membership consists of those persons so defined by C-VII,1 of the Constitution of the Association.

\subsection{Eligibility of New Members} \label{IV.2}
In addition to the requirements in \href{http://www.tbp.org/off/ConstBylaw.pdf}{C-II of the Constitution of the Association}, candidates must meet the following requirements:

\begin{enumerate}
    \item \subsecitem{a}{Residency}
    To be eligible for membership in the Chapter, an undergraduate student must have completed the equivalent of two full-time terms at the University of Michigan--Ann Arbor with at least one term in the College of Engineering, or be within one semester of graduation from the University of Michigan--Ann Arbor College of Engineering, or have been invited to join Tau Beta Pi at their previous school.
    \item \subsecitem{b}{Graduate Student Eligibility}
    To be eligible for election, graduate students must have additionally completed 11 credit hours of coursework as a graduate student.
The determination of percent completion of coursework as required by C-II of the Constitution of the Association is to be made by a student's primary advisor or program coordinator.
Graduate students who were eligible at the time of graduation from their undergraduate institution are eligible for election as alumni members.
    \item \subsecitem{c}{Eligible Undergraduate Curricula}
    Undergraduates in the following curricula in the College of Engineering are eligible for membership in the chapter:
    \begin{enumerate}
        \item Aerospace Engineering
        \item Biomedical Engineering
        \item Chemical Engineering
        \item Civil Engineering
        \item Civil and Environmental Engineering
        \item Climate and Space Sciences And Engineering
        \item Computer Engineering
        \item Computer Science
        \item Data Science
        \item Earth System Science and Engineering
        \item Electrical Engineering
        \item Engineering
        \item Engineering Physics
        \item Industrial and Operations Engineering
        \item Materials Science and Engineering
        \item Mechanical Engineering
        \item Naval Architecture and Marine Engineering
        \item Nuclear Engineering and Radiological Sciences
        \item Robotics Engineering 
    \end{enumerate}
    \item \subsecitem{d}{Eligible Graduate Curricula}
    Graduate students in any of the curricula listed in Bylaw~\rref{IV.2.c} or any of the following curricula are eligible for membership in the chapter:
    \begin{enumerate}
% Added F22 and successfully petitioned at 2022 convention
        \item Automotive Engineering
        \item Computer Science and Engineering
        \item Construction Engineering and Management
        \item Electrical Engineering: Systems
        \item Energy Systems Engineering
        \item Environmental Engineering
        \item Financial Engineering
        \item Global Automotive and Manufacturing Engineering
        \item Integrated Microsystems
        \item Macromolecular Science and Engineering
        \item Manufacturing Engineering
        \item Pharmaceutical Engineering
        \item Plastics Engineering
% Added F22 and successfully petitioned at 2022 convention
        \item Space Engineering
        \item Structural Engineering
    \end{enumerate}
\end{enumerate}

\subsection{Process} \label{IV.3}
The process for joining is referred to as the election process.
The requirements and process for joining consist of those requirements and items outlined in C-III of the Association, and in the chapter Bylaws.

% Added in October 2017 to comply with CCI's ridiculous requirements for constitutions. See Bylaws for MI-G process.
\subsection{Removal} \label{IV.4}
The process for removing a member from membership is defined in C-I,5 of the the Constitution of the Association.

\subsection{Support} \label{IV.5}
Upon joining the organization, all members agree not to undermine the purpose or mission of the Michigan Gamma chapter of Tau Beta Pi.